Choose linear coordinates so that $x_0,\ldots,x_r$ have no common zero on $X$ and $x_1/x_0,\ldots,x_r/x_0$ form a separating transcendence basis of $K(X)/k$. Such choices exist: avoiding $X$ is a nonempty open condition on a codimension-$(r+1)$ linear subspace, while independence of the corresponding differentials is a nonempty open condition by (4.7A) and (4.8A).

By the primitive element theorem, a $k$-linear combination of $x_{r+1}/x_0,\ldots,x_n/x_0$ generates $K(X)$ over $k(x_1/x_0,\ldots,x_r/x_0)$. Indeed, in a normal closure, failure to separate two embeddings imposes one of finitely many proper linear conditions on its coefficients, and $k$ is infinite. After another linear change, take this combination to be $x_{r+1}/x_0$.

Choose any $P\in Z(x_0,\ldots,x_{r+1})$, which is nonempty since $n\ge r+2$ and disjoint from $X$. Projection from $P$ to a complementary hyperplane is a morphism on $X$. Its coordinate linear forms include $x_0,\ldots,x_{r+1}$, so the function field of its image contains $k(x_1/x_0,\ldots,x_{r+1}/x_0)=K(X)$. Thus it induces an isomorphism of function fields and is birational onto its image.
